\documentclass[12pt,a4paper]{ctexart}

\usepackage[a4paper,margin=2.6cm]{geometry}
\usepackage{amsmath,amssymb,amsthm,mathtools,bm}
\usepackage{physics}
\usepackage{esint}
\usepackage{booktabs}
\usepackage{enumitem}
\usepackage[hypertexnames=false]{hyperref}

\hypersetup{
  colorlinks=true,
  linkcolor=blue,
  urlcolor=blue,
  citecolor=blue
}

\newtheorem{theorem}{定理}[section]
\newtheorem{definition}[theorem]{定义}
\newtheorem{proposition}[theorem]{命题}
\newtheorem{corollary}[theorem]{推论}
\theoremstyle{remark}
\newtheorem{remark}{注记}[section]
\newtheorem{example}{例}[section]

\title{数学分析复习提纲}
\author{基于教材《数学分析》第一册、第二册整理}
\date{}

\begin{document}

\maketitle

\begin{abstract}
本提纲系统整理数学分析课程的全部核心考点，涵盖集合与函数、极限与连续、
单变量微积分学、常微分方程、空间解析几何、多元函数微积分学及级数论等内容。
每一章包含核心定义、定理、方法和典型公式，按章—节—知识点三层结构组织。
定理编号与原始教材保持一致，符号系统严格遵循教材约定。
\end{abstract}

\tableofcontents

% ==================== 第一章 引论 ====================
\section{引论}

\subsection{集合}

\begin{definition}[集合与元素]
具有一定性质的事物的全体称为\textbf{集合}。构成集合的事物称为该集合的\textbf{元素}。
用 $a \in A$ 表示 $a$ 属于集合 $A$，用 $a \notin A$ 表示 $a$ 不属于集合 $A$。
\end{definition}

\begin{definition}[常用集合记号]
$\mathbb{N}$：自然数集；$\mathbb{Z}$：整数集；$\mathbb{Q}$：有理数集；
$\mathbb{R}$：实数集；$\mathbb{C}$：复数集；$\emptyset$：空集。
\end{definition}

\begin{definition}[集合的运算]
设 $A,B$ 为集合：
\begin{itemize}
  \item 并集：$A \cup B = \{x \mid x \in A \text{ 或 } x \in B\}$
  \item 交集：$A \cap B = \{x \mid x \in A \text{ 且 } x \in B\}$
  \item 差集：$A \setminus B = \{x \mid x \in A \text{ 且 } x \notin B\}$
  \item 补集：$A^c = \{x \mid x \notin A\}$（在全空间明确时使用）
\end{itemize}
\end{definition}

\begin{definition}[区间]
设 $a,b \in \mathbb{R}$ 且 $a < b$：
\begin{itemize}
  \item 闭区间：$[a,b] = \{x \mid a \le x \le b\}$
  \item 开区间：$(a,b) = \{x \mid a < x < b\}$
  \item 半开半闭区间：$[a,b)$，$(a,b]$
  \item 无穷区间：$[a,+\infty)$，$(-\infty,b]$ 等
\end{itemize}
\end{definition}

\begin{definition}[邻域]
设 $a \in \mathbb{R}$，$\delta > 0$，称 $(a-\delta, a+\delta)$ 为点 $a$ 的 $\delta$-\textbf{邻域}，
记作 $U(a,\delta)$。称 $(a-\delta, a) \cup (a, a+\delta)$ 为点 $a$ 的\textbf{去心邻域}，
记作 $\mathring{U}(a,\delta)$。
\end{definition}

\begin{remark}[确界原理]
非空有上界的实数集必有上确界；非空有下界的实数集必有下确界。
这是实数理论的基本定理之一，也是极限理论的基础。
\end{remark}

\subsection{函数}

\begin{definition}[函数]
设 $D \subset \mathbb{R}$ 为非空集合。若对每个 $x \in D$，按确定的法则 $f$，
有唯一的 $y \in \mathbb{R}$ 与之对应，则称 $f$ 是定义在 $D$ 上的\textbf{函数}，
记作 $y = f(x)$，$x \in D$。$D$ 称为定义域，$f(D) = \{f(x) \mid x \in D\}$ 称为值域。
\end{definition}

\begin{definition}[函数的性质]
\begin{itemize}
  \item \textbf{有界性}：若 $\exists M > 0$，$\forall x \in D$，有 $|f(x)| \le M$，则称 $f$ 在 $D$ 上有界。
  \item \textbf{单调性}：若 $\forall x_1,x_2 \in D$，$x_1 < x_2 \implies f(x_1) \le f(x_2)$（或 $\ge$），
    则称 $f$ 单调递增（或递减）。
  \item \textbf{奇偶性}：$f(-x) = -f(x)$ 为奇函数；$f(-x) = f(x)$ 为偶函数。
  \item \textbf{周期性}：若 $\exists T > 0$，使 $f(x+T) = f(x)$ 对所有 $x$ 成立，则 $T$ 为周期。
\end{itemize}
\end{definition}

\begin{definition}[基本初等函数]
常数函数、幂函数 $x^{\alpha}$、指数函数 $a^x$ $(a>0,a\neq 1)$、
对数函数 $\log_a x$ $(a>0,a\neq 1)$、三角函数、反三角函数。
\end{definition}

\begin{definition}[复合函数与反函数]
设 $y = f(u)$，$u = g(x)$，则 $y = f(g(x))$ 为 $f$ 与 $g$ 的\textbf{复合函数}。
若 $f$ 为单射，则存在\textbf{反函数} $f^{-1}$，满足 $f^{-1}(f(x)) = x$。
\end{definition}

% ==================== 第二章 极限与连续 ====================
\section{极限与连续}

\subsection{数列的极限}

\begin{definition}[数列极限]
设 $\{a_n\}$ 为数列，$a \in \mathbb{R}$。若 $\forall \varepsilon > 0$，
$\exists N \in \mathbb{N}$，当 $n > N$ 时恒有 $|a_n - a| < \varepsilon$，
则称数列 $\{a_n\}$ 收敛于 $a$，记作 $\lim\limits_{n\to\infty} a_n = a$。
\end{definition}

\begin{theorem}[收敛数列的性质]
收敛数列必为有界数列；极限若存在则唯一。
\end{theorem}

\begin{theorem}[极限的四则运算法则]
若 $\lim a_n = a$，$\lim b_n = b$，则
\begin{itemize}
  \item $\lim (a_n \pm b_n) = a \pm b$
  \item $\lim (a_n b_n) = ab$
  \item 若 $b \neq 0$，则 $\lim (a_n / b_n) = a/b$
\end{itemize}
\end{theorem}

\begin{theorem}[夹逼定理]
若 $\exists N$，使当 $n > N$ 时 $a_n \le b_n \le c_n$，且 $\lim a_n = \lim c_n = L$，
则 $\lim b_n = L$。
\end{theorem}

\begin{theorem}[单调有界定理]
单调递增（递减）且有上（下）界的数列必收敛。
\end{theorem}

\begin{definition}[重要极限]
$\displaystyle\lim_{n\to\infty}\left(1+\frac{1}{n}\right)^n = e$。
\end{definition}

\subsection{函数的极限}

\begin{definition}[函数极限（$x \to x_0$）]
设 $f$ 在 $x_0$ 的某去心邻域内有定义。若 $\forall \varepsilon > 0$，
$\exists \delta > 0$，当 $0 < |x - x_0| < \delta$ 时恒有 $|f(x) - A| < \varepsilon$，
则称 $x \to x_0$ 时 $f(x)$ 以 $A$ 为极限，记作 $\lim\limits_{x\to x_0} f(x) = A$。
\end{definition}

\begin{definition}[单侧极限]
左极限：$\lim\limits_{x\to x_0^-} f(x) = A$；右极限：$\lim\limits_{x\to x_0^+} f(x) = A$。
$\lim_{x\to x_0} f(x) = A$ 的充要条件是左右极限均存在且相等。
\end{definition}

\begin{definition}[函数极限（$x \to \infty$）]
$\lim\limits_{x\to\infty} f(x) = A$ 当且仅当
$\forall \varepsilon > 0$，$\exists X > 0$，当 $|x| > X$ 时 $|f(x)-A|<\varepsilon$。
\end{definition}

\begin{theorem}[函数极限的四则运算法则]
若 $\lim f(x) = A$，$\lim g(x) = B$，则和、差、积、商的极限法则与数列类似，
商的情形需 $B \neq 0$。
\end{theorem}

\begin{theorem}[复合函数的极限]
若 $\lim_{x\to x_0} g(x) = u_0$，$\lim_{u\to u_0} f(u) = A$，且在 $x_0$ 附近
$g(x) \neq u_0$（或 $f$ 在 $u_0$ 连续），则 $\lim_{x\to x_0} f(g(x)) = A$。
\end{theorem}

\begin{definition}[两个重要极限]
$\displaystyle\lim_{x\to 0}\frac{\sin x}{x}=1$；
$\displaystyle\lim_{x\to 0}(1+x)^{1/x}=e$。
\end{definition}

\subsection{无穷小与无穷大}

\begin{definition}[无穷小]
若 $\lim f(x) = 0$，则称 $f(x)$ 为该极限过程中的\textbf{无穷小量}。
\end{definition}

\begin{definition}[无穷小的比较]
设 $\alpha,\beta$ 均为无穷小，且 $\beta \neq 0$：
\begin{itemize}
  \item 若 $\lim \alpha/\beta = 0$，则称 $\alpha$ 是比 $\beta$ 高阶的无穷小，记作 $\alpha = o(\beta)$。
  \item 若 $\lim \alpha/\beta = c \neq 0$，则称 $\alpha$ 与 $\beta$ 为同阶无穷小。
  \item 若 $\lim \alpha/\beta = 1$，则称 $\alpha$ 与 $\beta$ 为等价无穷小，记作 $\alpha \sim \beta$。
\end{itemize}
\end{definition}

\begin{definition}[无穷大]
若 $\forall M > 0$，$\exists \delta > 0$（或 $X > 0$），使相应的 $x$ 满足 $|f(x)| > M$，
则称 $f(x)$ 为该过程中的\textbf{无穷大量}，记作 $\lim f(x) = \infty$。
\end{definition}

\subsection{连续函数}

\begin{definition}[连续性]
若 $\lim_{x\to x_0} f(x) = f(x_0)$，则称 $f$ 在 $x_0$ 处\textbf{连续}。
用 $\varepsilon$-$\delta$ 语言：$\forall \varepsilon > 0$，$\exists \delta > 0$，
当 $|x-x_0|<\delta$ 时 $|f(x)-f(x_0)|<\varepsilon$。
\end{definition}

\begin{definition}[间断点分类]
\begin{itemize}
  \item 第一类间断点：左右极限均存在（可去间断点、跳跃间断点）
  \item 第二类间断点：至少一侧极限不存在（无穷间断点、振荡间断点等）
\end{itemize}
\end{definition}

\begin{theorem}[连续函数的四则运算]
连续函数的和、差、积、商（分母不为零）仍为连续函数。
\end{theorem}

\begin{theorem}[复合函数的连续性]
若 $g$ 在 $x_0$ 连续，$f$ 在 $g(x_0)$ 连续，则 $f \circ g$ 在 $x_0$ 连续。
\end{theorem}

\begin{theorem}[反函数的连续性]
严格单调的连续函数其反函数也是严格单调的连续函数。
\end{theorem}

\subsection{闭区间上连续函数的性质}

\begin{theorem}[有界性定理]
闭区间上的连续函数必有界。
\end{theorem}

\begin{theorem}[最值定理]
闭区间上的连续函数必能取到最大值和最小值。
\end{theorem}

\begin{theorem}[介值定理]
若 $f$ 在 $[a,b]$ 上连续，且 $f(a) \neq f(b)$，则对于 $f(a)$ 与 $f(b)$
之间的任意值 $\mu$，存在 $\xi \in (a,b)$ 使 $f(\xi) = \mu$。
\end{theorem}

\begin{theorem}[零点存在定理]
若 $f$ 在 $[a,b]$ 上连续且 $f(a)f(b) < 0$，则存在 $\xi \in (a,b)$ 使 $f(\xi) = 0$。
\end{theorem}

\begin{theorem}[一致连续性]
若 $f$ 在闭区间 $[a,b]$ 上连续，则 $f$ 在 $[a,b]$ 上一致连续。
\end{theorem}

\begin{definition}[一致连续]
$f$ 在区间 $I$ 上一致连续，若 $\forall \varepsilon > 0$，$\exists \delta > 0$，
使对 $\forall x_1,x_2 \in I$，当 $|x_1-x_2|<\delta$ 时 $|f(x_1)-f(x_2)|<\varepsilon$。
注意与普通连续的区别：$\delta$ 不依赖于点 $x_0$ 的选取。
\end{definition}

% ==================== 第三章 单变量函数的微分学 ====================
\section{单变量函数的微分学}

\subsection{导数}

\begin{definition}[导数]
设函数 $f$ 在 $x_0$ 的某邻域内有定义。若极限
\[
f'(x_0) = \lim_{\Delta x \to 0} \frac{f(x_0+\Delta x)-f(x_0)}{\Delta x}
= \lim_{x \to x_0} \frac{f(x)-f(x_0)}{x-x_0}
\]
存在，则称 $f$ 在 $x_0$ 处\textbf{可导}，该极限值称为 $f$ 在 $x_0$ 的\textbf{导数}。
\end{definition}

\begin{theorem}[可导与连续]
可导必连续；连续不一定可导（如 $f(x)=|x|$ 在 $x=0$）。
\end{theorem}

\begin{definition}[高阶导数]
$f^{(n)}(x) = (f^{(n-1)}(x))'$，特别地 $f''(x)$ 为二阶导数。
\end{definition}

\subsection{微分}

\begin{definition}[微分]
若 $f$ 在 $x_0$ 可导，则 $f$ 在 $x_0$ 的\textbf{微分}为
$df = f'(x_0) \Delta x = f'(x_0) dx$。
\end{definition}

\begin{remark}
微分是函数改变量 $\Delta y = f(x_0+\Delta x)-f(x_0)$ 的线性主部：
$\Delta y = f'(x_0) \Delta x + o(\Delta x)$。
一阶微分形式不变性：无论 $x$ 是自变量还是中间变量，$dy = f'(x)dx$ 形式不变。
\end{remark}

\subsection{求导法则与高阶导数}

\begin{theorem}[基本求导法则]
\begin{itemize}
  \item $(u \pm v)' = u' \pm v'$
  \item $(uv)' = u'v + uv'$
  \item $(\frac{u}{v})' = \frac{u'v - uv'}{v^2}$（$v \neq 0$）
  \item 链式法则：$(f(g(x)))' = f'(g(x)) \cdot g'(x)$
  \item 反函数求导：$(f^{-1})'(y) = 1/f'(x)$，其中 $y = f(x)$
\end{itemize}
\end{theorem}

\begin{definition}[基本初等函数的导数公式]
$(x^\alpha)' = \alpha x^{\alpha-1}$；$(a^x)' = a^x\ln a$；$(\ln x)' = 1/x$；
$(\sin x)' = \cos x$；$(\cos x)' = -\sin x$；$(\tan x)' = \sec^2 x$；等等。
\end{definition}

\begin{theorem}[Leibniz 公式]
$(uv)^{(n)} = \sum_{k=0}^n \binom{n}{k} u^{(n-k)} v^{(k)}$。
\end{theorem}

\begin{definition}[参数方程求导]
若 $x = \varphi(t)$，$y = \psi(t)$，则 $\frac{dy}{dx} = \frac{\psi'(t)}{\varphi'(t)}$，
$\frac{d^2y}{dx^2} = \frac{\psi''(t)\varphi'(t)-\psi'(t)\varphi''(t)}{[\varphi'(t)]^3}$。
\end{definition}

\subsection{微分学中值定理}

\begin{theorem}[Rolle 定理]
若 $f$ 在 $[a,b]$ 上连续，在 $(a,b)$ 内可导，且 $f(a)=f(b)$，则存在
$\xi \in (a,b)$ 使 $f'(\xi) = 0$。
\end{theorem}

\begin{theorem}[Lagrange 中值定理]
若 $f$ 在 $[a,b]$ 上连续，在 $(a,b)$ 内可导，则存在 $\xi \in (a,b)$ 使
\[
f'(\xi) = \frac{f(b)-f(a)}{b-a}.
\]
\end{theorem}

\begin{corollary}[导数与单调性]
若在区间 $I$ 上 $f'(x) \equiv 0$，则 $f$ 在 $I$ 上为常数。
若 $f'(x) \ge 0$（$\le 0$），则 $f$ 单调递增（递减）。
\end{corollary}

\begin{theorem}[Cauchy 中值定理]
若 $f,g$ 在 $[a,b]$ 上连续，在 $(a,b)$ 内可导，且 $g'(x) \neq 0$，则存在 $\xi \in (a,b)$ 使
\[
\frac{f'(\xi)}{g'(\xi)} = \frac{f(b)-f(a)}{g(b)-g(a)}.
\]
\end{theorem}

\begin{theorem}[L'H\^opital 法则]
若 $\lim f(x) = \lim g(x) = 0$（或 $\infty$），且 $\lim f'(x)/g'(x)$ 存在，
则 $\lim f(x)/g(x) = \lim f'(x)/g'(x)$。适用于 $\frac{0}{0}$ 型和 $\frac{\infty}{\infty}$ 型。
\end{theorem}

\subsection{Taylor 公式}

\begin{theorem}[带 Lagrange 余项的 Taylor 公式]
若 $f$ 在 $x_0$ 的某邻域内有 $n+1$ 阶导数，则
\[
f(x) = \sum_{k=0}^n \frac{f^{(k)}(x_0)}{k!}(x-x_0)^k + R_n(x),
\]
其中 $R_n(x) = \frac{f^{(n+1)}(\xi)}{(n+1)!}(x-x_0)^{n+1}$，$\xi$ 在 $x_0$ 与 $x$ 之间。
\end{theorem}

\begin{theorem}[带 Peano 余项的 Taylor 公式]
若 $f$ 在 $x_0$ 处 $n$ 阶可导，则
\[
f(x) = \sum_{k=0}^n \frac{f^{(k)}(x_0)}{k!}(x-x_0)^k + o((x-x_0)^n).
\]
\end{theorem}

\begin{definition}[Maclaurin 公式（$x_0=0$）]
$e^x = 1 + x + \frac{x^2}{2!} + \cdots + \frac{x^n}{n!} + o(x^n)$；
$\sin x = x - \frac{x^3}{3!} + \frac{x^5}{5!} - \cdots$；
$\cos x = 1 - \frac{x^2}{2!} + \frac{x^4}{4!} - \cdots$；
$\ln(1+x) = x - \frac{x^2}{2} + \frac{x^3}{3} - \cdots$；
$(1+x)^\alpha = 1 + \alpha x + \frac{\alpha(\alpha-1)}{2!}x^2 + \cdots$。
\end{definition}

\subsection{应用}

\begin{definition}[极值]
若在 $x_0$ 的某邻域内恒有 $f(x) \le f(x_0)$（或 $\ge$），则 $f(x_0)$ 为极大值（极小值）。
\end{definition}

\begin{theorem}[极值的必要条件]
若 $f$ 在 $x_0$ 可导且取极值，则 $f'(x_0) = 0$（驻点）。
\end{theorem}

\begin{theorem}[极值的充分条件]
(1) 若 $f'$ 在 $x_0$ 左右变号，则 $x_0$ 为极值点；(2) 若 $f'(x_0)=0$ 且 $f''(x_0) \neq 0$，
则当 $f''(x_0) > 0$ 时为极小值，$f''(x_0) < 0$ 时为极大值。
\end{theorem}

\begin{definition}[凸性与拐点]
若 $f''(x) > 0$，$f$ 是下凸（凹向上）的；$f''(x) < 0$，$f$ 是上凸（凹向下）的。
$f''$ 变号的点称为拐点。
\end{definition}

\begin{definition}[渐近线]
水平渐近线：$\lim_{x\to\infty} f(x) = b$；垂直渐近线：$\lim_{x\to a} f(x) = \infty$；
斜渐近线：$y = kx + b$，其中 $k = \lim_{x\to\infty} f(x)/x$，$b = \lim_{x\to\infty}(f(x)-kx)$。
\end{definition}

% ==================== 第四章 单变量函数的积分学 ====================
\section{单变量函数的积分学}

\subsection{不定积分}

\begin{definition}[原函数与不定积分]
若 $F'(x) = f(x)$，则称 $F$ 为 $f$ 的一个\textbf{原函数}。
$f$ 的不定积分为 $\int f(x) dx = F(x) + C$，其中 $C$ 为任意常数。
\end{definition}

\begin{definition}[基本积分公式]
$\int x^\alpha dx = \frac{x^{\alpha+1}}{\alpha+1}+C\ (\alpha\neq -1)$；
$\int \frac{1}{x} dx = \ln|x| + C$；
$\int e^x dx = e^x + C$；
$\int \sin x dx = -\cos x + C$；
$\int \cos x dx = \sin x + C$；等等。
\end{definition}

\subsection{积分法}

\begin{definition}[换元积分法]
第一类换元（凑微分）：$\int f(g(x))g'(x)dx = \int f(u) du$。
第二类换元（变量代换）：$\int f(x)dx = \int f(g(t))g'(t)dt$。
\end{definition}

\begin{theorem}[分部积分法]
$\int u dv = uv - \int v du$，即 $\int u v' dx = uv - \int u' v dx$。
\end{theorem}

\begin{definition}[有理函数的积分]
将有理函数分解为部分分式之和，分别积分。分母因式分解后，对应形式有
$\frac{A}{x-a}$，$\frac{A}{(x-a)^k}$，$\frac{Ax+B}{x^2+px+q}$，$\frac{Ax+B}{(x^2+px+q)^k}$ 等。
\end{definition}

\begin{definition}[三角函数有理式的积分]
令 $t = \tan\frac{x}{2}$（万能代换），则 $\sin x = \frac{2t}{1+t^2}$，
$\cos x = \frac{1-t^2}{1+t^2}$，$dx = \frac{2}{1+t^2}dt$。
\end{definition}

\begin{definition}[简单无理函数的积分]
含 $\sqrt{ax^2+bx+c}$ 的被积函数可通过三角代换或 Euler 代换处理。
\end{definition}

\subsection{定积分}

\begin{definition}[Riemann 积分]
设 $f$ 在 $[a,b]$ 上有定义。对 $[a,b]$ 作分割
$a = x_0 < x_1 < \cdots < x_n = b$，
取 $\xi_i \in [x_{i-1}, x_i]$，作 Riemann 和
$\sum_{i=1}^n f(\xi_i)\Delta x_i$。
若当 $\max \Delta x_i \to 0$ 时该和的极限存在且与分割和取点无关，
则称 $f$ 在 $[a,b]$ 上 Riemann 可积，记积分值为
$\int_a^b f(x) dx$。
\end{definition}

\begin{theorem}[可积的充分条件]
闭区间上的连续函数可积；闭区间上只有有限个间断点的有界函数可积；
闭区间上的单调函数可积。
\end{theorem}

\begin{theorem}[Newton-Leibniz 公式]
若 $f$ 在 $[a,b]$ 上可积，$F$ 是 $f$ 的一个原函数，则
\[
\int_a^b f(x) dx = F(b) - F(a) = \left.F(x)\right|_a^b.
\]
\end{theorem}

\begin{theorem}[定积分的性质]
\begin{itemize}
  \item 线性性：$\int_a^b (\alpha f + \beta g) = \alpha \int_a^b f + \beta \int_a^b g$
  \item 区间可加性：$\int_a^b f = \int_a^c f + \int_c^b f$
  \item 保序性：若 $f \le g$，则 $\int_a^b f \le \int_a^b g$
  \item 绝对值不等式：$|\int_a^b f| \le \int_a^b |f|$
\end{itemize}
\end{theorem}

\begin{theorem}[积分中值定理]
若 $f$ 在 $[a,b]$ 上连续，则存在 $\xi \in [a,b]$ 使
$\int_a^b f(x) dx = f(\xi)(b-a)$。
\end{theorem}

\begin{theorem}[原函数存在定理（微积分基本定理）]
若 $f$ 在 $[a,b]$ 上连续，则变上限积分 $\Phi(x) = \int_a^x f(t) dt$ 在 $[a,b]$
上可导，且 $\Phi'(x) = f(x)$。
\end{theorem}

\begin{theorem}[定积分的换元与分部积分]
换元：$\int_a^b f(x) dx = \int_{\alpha}^{\beta} f(g(t)) g'(t) dt$，其中 $x = g(t)$。
分部：$\int_a^b u dv = uv|_a^b - \int_a^b v du$。
\end{theorem}

\subsection{广义积分}

\begin{definition}[无穷区间上的广义积分]
$\int_a^{+\infty} f(x) dx = \lim_{b\to+\infty} \int_a^b f(x) dx$。
若极限存在，称广义积分收敛，否则发散。
\end{definition}

\begin{definition}[无界函数的广义积分（瑕积分）]
若 $f$ 在 $b$ 点无界，则 $\int_a^b f(x) dx = \lim_{\varepsilon\to 0^+} \int_a^{b-\varepsilon} f(x) dx$。
\end{definition}

\begin{theorem}[比较判别法]
若 $0 \le f(x) \le g(x)$，则 $\int_a^{+\infty} g(x)dx$ 收敛 $\implies$
$\int_a^{+\infty} f(x)dx$ 收敛；$\int_a^{+\infty} f(x)dx$ 发散 $\implies$
$\int_a^{+\infty} g(x)dx$ 发散。
\end{theorem}

\begin{theorem}[绝对收敛与条件收敛]
若 $\int_a^{+\infty} |f(x)| dx$ 收敛，则 $\int_a^{+\infty} f(x) dx$ 收敛（绝对收敛）。
若 $\int_a^{+\infty} f(x) dx$ 收敛而 $\int_a^{+\infty} |f(x)| dx$ 发散，则为条件收敛。
\end{theorem}

\begin{definition}[$\Gamma$ 函数与 $\mathrm{B}$ 函数]
$\Gamma(s) = \int_0^{+\infty} x^{s-1} e^{-x} dx\ (s > 0)$；
$\mathrm{B}(p,q) = \int_0^1 x^{p-1}(1-x)^{q-1} dx\ (p,q > 0)$。
$\mathrm{B}(p,q) = \frac{\Gamma(p)\Gamma(q)}{\Gamma(p+q)}$。
\end{definition}

\subsection{应用}

定积分的应用包括：平面图形的面积、旋转体的体积、平面曲线的弧长、
旋转曲面的面积、物理应用（功、压力、质心等）。

\begin{definition}[面积公式]
极坐标下 $S = \frac{1}{2}\int_{\alpha}^{\beta} r^2(\theta) d\theta$。
参数方程下 $S = \int_a^b y(t) x'(t) dt$。
\end{definition}

\begin{definition}[弧长公式]
$s = \int_a^b \sqrt{1+(y')^2} dx$；参数方程 $s = \int_{t_1}^{t_2} \sqrt{(x')^2+(y')^2} dt$。
\end{definition}

\begin{definition}[旋转体体积]
绕 $x$ 轴：$V = \pi \int_a^b y^2 dx$；绕 $y$ 轴：$V = 2\pi \int_a^b xy dx$（柱壳法）。
\end{definition}

% ==================== 第五章 常微分方程 ====================
\section{常微分方程}

\subsection{基本概念}

\begin{definition}[微分方程]
含有未知函数及其导数的方程。未知函数为一元函数的称为\textbf{常微分方程}。
方程中最高阶导数的阶数称为方程的\textbf{阶}。
\end{definition}

\begin{definition}[通解与特解]
$n$ 阶方程的通解含有 $n$ 个独立的任意常数。确定了常数值的解为特解。
定解条件（初值条件）确定了特解。
\end{definition}

\subsection{初等积分法}

\begin{definition}[可分离变量方程]
$\frac{dy}{dx} = f(x)g(y)$，分离变量后两边积分：
$\int \frac{dy}{g(y)} = \int f(x) dx$。
\end{definition}

\begin{definition}[齐次方程]
$\frac{dy}{dx} = \varphi(\frac{y}{x})$。令 $u = y/x$，则 $y = ux$，
$y' = u + x u'$，化为可分离变量方程。
\end{definition}

\begin{definition}[一阶线性方程]
$y' + P(x)y = Q(x)$。通解公式：
\[
y = e^{-\int P dx}\left(C + \int Q e^{\int P dx} dx\right).
\]
\end{definition}

\begin{definition}[Bernoulli 方程]
$y' + P(x)y = Q(x)y^n\ (n \neq 0,1)$。令 $z = y^{1-n}$ 化为一阶线性方程。
\end{definition}

\begin{definition}[全微分方程]
$P(x,y)dx + Q(x,y)dy = 0$，若 $\frac{\partial P}{\partial y} = \frac{\partial Q}{\partial x}$，
则为全微分方程，存在 $u(x,y)$ 使 $du = Pdx + Qdy$，通解为 $u(x,y) = C$。
\end{definition}

\subsection{线性微分方程组}

\begin{definition}[一阶线性微分方程组]
$\frac{d\mathbf{y}}{dx} = \mathbf{A}(x)\mathbf{y} + \mathbf{f}(x)$，
其中 $\mathbf{y} = (y_1,\dots,y_n)^T$，$\mathbf{A}(x)$ 为 $n\times n$ 矩阵。
\end{definition}

\begin{theorem}[齐次方程组解的结构]
齐次方程 $\mathbf{y}' = \mathbf{A}(x)\mathbf{y}$ 的解构成 $n$ 维线性空间。
$n$ 个线性无关的特解构成基本解组（基本解矩阵）。
\end{theorem}

\begin{theorem}[常数变易法]
非齐次方程的特解可通过常数变易法求得：
$\mathbf{y}_p = \mathbf{Y}(x)\int \mathbf{Y}^{-1}(x)\mathbf{f}(x)dx$，
其中 $\mathbf{Y}(x)$ 为基本解矩阵。
\end{theorem}

\subsection{高阶线性方程}

\begin{definition}[$n$ 阶线性微分方程]
$y^{(n)} + a_1(x)y^{(n-1)} + \cdots + a_n(x)y = f(x)$。
当 $f(x) \equiv 0$ 时为齐次，否则为非齐次。
\end{definition}

\begin{theorem}[齐次方程解的结构]
齐次线性方程的通解为 $n$ 个线性无关特解的线性组合：
$y = C_1 y_1 + C_2 y_2 + \cdots + C_n y_n$。
\end{theorem}

\begin{theorem}[常系数齐次线性方程]
$y^{(n)} + a_1 y^{(n-1)} + \cdots + a_n y = 0$，其中 $a_i$ 为常数。
其特征方程为 $\lambda^n + a_1\lambda^{n-1} + \cdots + a_n = 0$。
根据特征根（单实根、重实根、共轭复根）写出对应的基解：
\begin{itemize}
  \item 单实根 $\lambda$：$e^{\lambda x}$
  \item $k$ 重实根 $\lambda$：$e^{\lambda x}, xe^{\lambda x}, \dots, x^{k-1}e^{\lambda x}$
  \item 共轭复根 $\alpha \pm i\beta$：$e^{\alpha x}\cos\beta x, e^{\alpha x}\sin\beta x$
\end{itemize}
\end{theorem}

\begin{theorem}[非齐次方程的待定系数法]
对于 $f(x) = P_m(x)e^{\lambda x}$ 或 $f(x) = e^{\alpha x}(P_m(x)\cos\beta x + Q_n(x)\sin\beta x)$
的形式，可设相应形式的特解待定系数求解。
\end{theorem}

% ==================== 第六章 空间解析几何 ====================
\section{空间解析几何}

\subsection{向量代数}

\begin{definition}[向量]
既有大小又有方向的量。向量 $\mathbf{a}$ 的模记作 $|\mathbf{a}|$。
零向量 $\mathbf{0}$ 模为零。
\end{definition}

\begin{definition}[向量的运算]
\begin{itemize}
  \item 加法：平行四边形法则或三角形法则
  \item 数乘：$k\mathbf{a}$ 与 $\mathbf{a}$ 共线，$|k\mathbf{a}| = |k||\mathbf{a}|$
  \item 内积（点积）：$\mathbf{a} \cdot \mathbf{b} = |\mathbf{a}||\mathbf{b}|\cos\theta$
  \item 外积（叉积）：$\mathbf{a} \times \mathbf{b}$ 垂直于 $\mathbf{a},\mathbf{b}$ 所在平面，
    $|\mathbf{a} \times \mathbf{b}| = |\mathbf{a}||\mathbf{b}|\sin\theta$
  \item 混合积：$(\mathbf{a},\mathbf{b},\mathbf{c}) = (\mathbf{a} \times \mathbf{b}) \cdot \mathbf{c}$
\end{itemize}
\end{definition}

\begin{definition}[向量共线与共面]
$\mathbf{a} \parallel \mathbf{b} \iff \mathbf{a} \times \mathbf{b} = \mathbf{0}$；
$\mathbf{a},\mathbf{b},\mathbf{c}$ 共面 $\iff (\mathbf{a},\mathbf{b},\mathbf{c}) = 0$。
\end{definition}

\subsection{平面与直线}

\begin{definition}[平面方程]
\begin{itemize}
  \item 点法式：$A(x-x_0)+B(y-y_0)+C(z-z_0)=0$，法向量 $\mathbf{n}=(A,B,C)$
  \item 一般式：$Ax+By+Cz+D=0$
  \item 截距式：$\frac{x}{a}+\frac{y}{b}+\frac{z}{c}=1$
\end{itemize}
\end{definition}

\begin{definition}[直线方程]
\begin{itemize}
  \item 点向式：$\frac{x-x_0}{l}=\frac{y-y_0}{m}=\frac{z-z_0}{n}$，方向向量 $\mathbf{s}=(l,m,n)$
  \item 参数式：$x=x_0+lt$，$y=y_0+mt$，$z=z_0+nt$
  \item 一般式：$\begin{cases} A_1x+B_1y+C_1z+D_1=0 \\ A_2x+B_2y+C_2z+D_2=0 \end{cases}$
\end{itemize}
\end{definition}

\begin{definition}[点/线/面之间的位置关系]
点到平面的距离：$d = \frac{|Ax_0+By_0+Cz_0+D|}{\sqrt{A^2+B^2+C^2}}$。
两平面的夹角即法向量的夹角；直线与平面的夹角为方向向量与法向量夹角的余角。
\end{definition}

\subsection{曲面与空间曲线}

\begin{definition}[常见二次曲面]
\begin{itemize}
  \item 球面：$(x-a)^2+(y-b)^2+(z-c)^2=R^2$
  \item 椭球面：$\frac{x^2}{a^2}+\frac{y^2}{b^2}+\frac{z^2}{c^2}=1$
  \item 单叶双曲面：$\frac{x^2}{a^2}+\frac{y^2}{b^2}-\frac{z^2}{c^2}=1$
  \item 双叶双曲面：$\frac{x^2}{a^2}+\frac{y^2}{b^2}-\frac{z^2}{c^2}=-1$
  \item 椭圆抛物面：$z = \frac{x^2}{a^2}+\frac{y^2}{b^2}$
  \item 双曲抛物面（马鞍面）：$z = \frac{x^2}{a^2}-\frac{y^2}{b^2}$
  \item 二次锥面：$\frac{x^2}{a^2}+\frac{y^2}{b^2}-\frac{z^2}{c^2}=0$
\end{itemize}
\end{definition}

% ==================== 第七章 多元函数微分学 ====================
\section{多元函数微分学}

\subsection{多元函数的极限与连续}

\begin{definition}[多元函数]
设 $D \subset \mathbb{R}^n$，映射 $f: D \to \mathbb{R}$ 称为 $n$ 元函数，
记作 $u = f(x_1,\dots,x_n)$ 或 $u = f(\mathbf{x})$。
\end{definition}

\begin{definition}[二重极限]
设 $f$ 在 $(x_0,y_0)$ 的某去心邻域内有定义。若 $\forall \varepsilon > 0$，
$\exists \delta > 0$，当 $0 < \sqrt{(x-x_0)^2+(y-y_0)^2} < \delta$ 时
$|f(x,y)-A|<\varepsilon$，则称 $\lim_{(x,y)\to(x_0,y_0)} f(x,y) = A$。
\end{definition}

\begin{remark}
二重极限存在要求沿任何路径趋于 $(x_0,y_0)$ 时极限均相同且等于 $A$。
若沿两条不同路径极限不同，则二重极限不存在。
\end{remark}

\begin{definition}[累次极限]
$\lim_{x\to x_0}\lim_{y\to y_0} f(x,y)$ 和 $\lim_{y\to y_0}\lim_{x\to x_0} f(x,y)$
称为累次极限。二重极限与累次极限不一定相等。
\end{definition}

\begin{definition}[连续性]
若 $\lim_{(x,y)\to(x_0,y_0)} f(x,y) = f(x_0,y_0)$，则 $f$ 在 $(x_0,y_0)$ 连续。
\end{definition}

\subsection{偏导数与全微分}

\begin{definition}[偏导数]
$f_x(x_0,y_0) = \lim_{\Delta x\to 0} \frac{f(x_0+\Delta x, y_0)-f(x_0,y_0)}{\Delta x}$。
$f_y$ 类似定义。
高阶偏导数：$f_{xx} = \frac{\partial^2 f}{\partial x^2}$，
$f_{xy} = \frac{\partial^2 f}{\partial x \partial y}$。
\end{definition}

\begin{theorem}[混合偏导数的对称性]
若 $f_{xy}$ 和 $f_{yx}$ 在 $(x_0,y_0)$ 连续，则 $f_{xy}(x_0,y_0) = f_{yx}(x_0,y_0)$。
\end{theorem}

\begin{definition}[全微分]
若 $f$ 在 $(x_0,y_0)$ 的全改变量可表示为
$\Delta z = A\Delta x + B\Delta y + o(\sqrt{\Delta x^2 + \Delta y^2})$，
则称 $f$ 在 $(x_0,y_0)$ 可微，$dz = A dx + B dy$。
此时 $A = f_x(x_0,y_0)$，$B = f_y(x_0,y_0)$。
\end{definition}

\begin{theorem}[可微的充分条件]
若 $f_x, f_y$ 在 $(x_0,y_0)$ 连续，则 $f$ 在 $(x_0,y_0)$ 可微。
\end{theorem}

\begin{theorem}[链式法则]
若 $z = f(u,v)$，$u = u(x,y)$，$v = v(x,y)$，则
$\frac{\partial z}{\partial x} = \frac{\partial f}{\partial u}\frac{\partial u}{\partial x}
+ \frac{\partial f}{\partial v}\frac{\partial v}{\partial x}$。
\end{theorem}

\subsection{方向导数与梯度}

\begin{definition}[方向导数]
$\frac{\partial f}{\partial \mathbf{l}} = \lim_{t\to 0^+} \frac{f(\mathbf{x}_0+t\mathbf{l})-f(\mathbf{x}_0)}{t}$，
其中 $\mathbf{l}$ 为单位向量。
若 $f$ 可微，则 $\frac{\partial f}{\partial \mathbf{l}} = \nabla f \cdot \mathbf{l}$。
\end{definition}

\begin{definition}[梯度]
$\nabla f = \operatorname{grad} f = (f_x, f_y, f_z)$。梯度方向是函数增长最快的方向。
\end{definition}

\subsection{多元函数的 Taylor 公式与极值}

\begin{theorem}[多元 Taylor 公式]
$f(\mathbf{x}) = f(\mathbf{x}_0) + \nabla f(\mathbf{x}_0) \cdot (\mathbf{x}-\mathbf{x}_0)
+ \frac{1}{2}(\mathbf{x}-\mathbf{x}_0)^T H_f(\mathbf{x}_0)(\mathbf{x}-\mathbf{x}_0) + \cdots$
其中 $H_f$ 为 Hessian 矩阵，$H_{ij} = \frac{\partial^2 f}{\partial x_i \partial x_j}$。
\end{theorem}

\begin{definition}[极值]
若在 $\mathbf{x}_0$ 的某邻域内 $f(\mathbf{x}) \le f(\mathbf{x}_0)$，则为极大值。
\end{definition}

\begin{theorem}[极值的必要条件]
若 $f$ 在 $\mathbf{x}_0$ 可微且取极值，则 $\nabla f(\mathbf{x}_0) = \mathbf{0}$（驻点）。
\end{theorem}

\begin{theorem}[极值的充分条件]
设 $\nabla f(\mathbf{x}_0) = \mathbf{0}$，考察 Hessian 矩阵 $H_f(\mathbf{x}_0)$：
\begin{itemize}
  \item $H$ 正定 $\implies$ 极小值
  \item $H$ 负定 $\implies$ 极大值
  \item $H$ 不定 $\implies$ 非极值（鞍点）
\end{itemize}
对于二元函数，可用判别式 $\Delta = f_{xx}f_{yy} - f_{xy}^2$：
若 $\Delta > 0$ 且 $f_{xx} > 0$ 则为极小，$f_{xx} < 0$ 则为极大；
若 $\Delta < 0$ 则非极值。
\end{theorem}

\begin{definition}[条件极值（Lagrange 乘数法）]
求 $f(\mathbf{x})$ 在条件 $g(\mathbf{x})=0$ 下的极值。构造 Lagrange 函数
$L(\mathbf{x},\lambda) = f(\mathbf{x}) + \lambda g(\mathbf{x})$，
解 $\nabla L = \mathbf{0}$。
多约束时：$L = f + \sum \lambda_i g_i$。
\end{definition}

% ==================== 第八章 多元函数积分学 ====================
\section{多元函数积分学}

\subsection{二重积分}

\begin{definition}[二重积分]
$\iint_D f(x,y) d\sigma = \lim_{\max \Delta\sigma_i \to 0} \sum f(\xi_i,\eta_i)\Delta\sigma_i$，
其中 $D$ 为有界闭区域。连续函数在可求积的闭区域上二重积分存在。
\end{definition}

\begin{theorem}[二重积分的计算]
直角坐标下：$\iint_D f(x,y) dxdy = \int_a^b dx \int_{\varphi_1(x)}^{\varphi_2(x)} f(x,y) dy$
（或先 $x$ 后 $y$）。
极坐标下：$\iint_D f(x,y) dxdy = \int_{\alpha}^{\beta} d\theta \int_{r_1(\theta)}^{r_2(\theta)} f(r\cos\theta,r\sin\theta) r dr$。
注意雅可比因子 $r$。
\end{theorem}

\begin{theorem}[二重积分的变量代换]
$\iint_D f(x,y) dxdy = \iint_{D'} f(x(u,v), y(u,v)) |J| dudv$，
其中 $J = \frac{\partial(x,y)}{\partial(u,v)}$ 为雅可比行列式。
\end{theorem}

\subsection{三重积分}

\begin{definition}[三重积分]
$\iiint_\Omega f(x,y,z) dV$。计算方法包括：直角坐标（截面法、投影法）、
柱坐标（$x=r\cos\theta,y=r\sin\theta,z=z$，$J=r$）、
球坐标（$x=\rho\sin\varphi\cos\theta,y=\rho\sin\varphi\sin\theta,z=\rho\cos\varphi$，
$J=\rho^2\sin\varphi$）。
\end{definition}

\subsection{曲线积分}

\begin{definition}[第一类曲线积分（对弧长的积分）]
$\int_L f(x,y) ds = \lim_{\max \Delta s_i \to 0} \sum f(\xi_i,\eta_i) \Delta s_i$。
计算方法：若 $L: y=y(x), a\le x\le b$，
则 $ds = \sqrt{1+(y')^2}dx$，$\int_L f ds = \int_a^b f(x,y(x))\sqrt{1+(y')^2}dx$。
参数方程下 $ds = \sqrt{(x'(t))^2+(y'(t))^2} dt$。
\end{definition}

\begin{definition}[第二类曲线积分（对坐标的积分）]
$\int_L Pdx + Qdy = \lim_{\max\Delta x_i\to 0} \sum [P(\xi_i,\eta_i)\Delta x_i + Q(\xi_i,\eta_i)\Delta y_i]$。
参数方程下：$\int_L Pdx+Qdy = \int_{t_1}^{t_2} [P x'(t) + Q y'(t)] dt$。
\end{definition}

\begin{theorem}[Green 公式]
设 $D$ 为由分段光滑的简单闭曲线 $L$ 围成的有界闭区域，$P,Q$ 在 $D$ 上有连续一阶偏导数，则
\[
\oint_L Pdx + Qdy = \iint_D \left(\frac{\partial Q}{\partial x} - \frac{\partial P}{\partial y}\right) dxdy.
\]
其中 $L$ 取正向（逆时针）。
\end{theorem}

\begin{theorem}[曲线积分与路径无关]
以下条件等价：在单连通区域 $D$ 内，
\begin{enumerate}
  \item 对 $D$ 内任意闭曲线 $L$，$\oint_L Pdx+Qdy=0$
  \item 曲线积分与路径无关
  \item $Pdx+Qdy$ 是某函数 $u(x,y)$ 的全微分
  \item $\frac{\partial P}{\partial y} = \frac{\partial Q}{\partial x}$ 在 $D$ 内处处成立
\end{enumerate}
\end{theorem}

\subsection{曲面积分}

\begin{definition}[第一类曲面积分（对面积的积分）]
$\iint_S f(x,y,z) dS$。计算方法：若 $S: z = z(x,y)$，$dS = \sqrt{1+z_x^2+z_y^2}dxdy$。
\end{definition}

\begin{definition}[第二类曲面积分（对坐标的积分）]
$\iint_S P dydz + Q dzdx + R dxdy$。正侧根据曲面定向确定符号。
\end{definition}

\begin{theorem}[Gauss 公式（散度定理）]
设 $\Omega$ 为由分片光滑闭曲面 $S$ 围成的有界闭区域，则
\[
\oiint_S P dydz + Q dzdx + R dxdy = \iiint_\Omega
\left(\frac{\partial P}{\partial x} + \frac{\partial Q}{\partial y} + \frac{\partial R}{\partial z}\right) dV.
\]
其中 $S$ 取外侧。记散度为 $\operatorname{div} \mathbf{F} = \frac{\partial P}{\partial x} + \frac{\partial Q}{\partial y} + \frac{\partial R}{\partial z}$。
\end{theorem}

\begin{theorem}[Stokes 公式]
设 $S$ 为以分段光滑闭曲线 $L$ 为边界的有向曲面，则
\[
\oint_L Pdx+Qdy+Rdz = \iint_S
\begin{vmatrix}
dydz & dzdx & dxdy \\
\frac{\partial}{\partial x} & \frac{\partial}{\partial y} & \frac{\partial}{\partial z} \\
P & Q & R
\end{vmatrix}.
\]
记旋度为 $\operatorname{rot} \mathbf{F} = \nabla \times \mathbf{F}$。
\end{theorem}

% ==================== 第九章 级数论 ====================
\section{级数论}

\subsection{数项级数}

\begin{definition}[数项级数]
$\sum_{n=1}^{\infty} a_n = \lim_{n\to\infty} S_n$，其中 $S_n = \sum_{k=1}^n a_k$ 为部分和。
若 $\{S_n\}$ 收敛，则级数收敛；否则发散。
\end{definition}

\begin{theorem}[级数收敛的必要条件]
若 $\sum a_n$ 收敛，则 $\lim a_n = 0$。其逆不真（如调和级数）。
\end{theorem}

\begin{definition}[正项级数及其判别法]
\begin{itemize}
  \item \textbf{比较判别法}：若 $0 \le a_n \le b_n$，则 $\sum b_n$ 收敛 $\implies$ $\sum a_n$ 收敛
  \item \textbf{比值判别法（d'Alembert）}：$\lim \frac{a_{n+1}}{a_n} = \rho$，
    $\rho<1$ 收敛，$\rho>1$ 发散，$\rho=1$ 不确定
  \item \textbf{根值判别法（Cauchy）}：$\lim \sqrt[n]{a_n} = \rho$，同上
  \item \textbf{积分判别法}：若 $f(x)$ 单调递减非负，则 $\sum f(n)$ 与 $\int_1^{\infty} f(x)dx$ 同敛散
  \item \textbf{$p$-级数}：$\sum 1/n^p$，$p>1$ 收敛，$p\le 1$ 发散
\end{itemize}
\end{definition}

\begin{definition}[交错级数]
$\sum (-1)^{n-1} a_n$，其中 $a_n > 0$。Leibniz 判别法：若 $a_n$ 单调递减趋于零，
则交错级数收敛。
\end{definition}

\begin{definition}[绝对收敛与条件收敛]
若 $\sum |a_n|$ 收敛，则 $\sum a_n$ 绝对收敛；
若 $\sum a_n$ 收敛但 $\sum |a_n|$ 发散，则为条件收敛。
绝对收敛级数可任意重排而不改变收敛性和和值。
\end{definition}

\subsection{函数项级数}

\begin{definition}[函数项级数、收敛域]
$\sum u_n(x)$。使级数收敛的 $x$ 全体称为收敛域。
\end{definition}

\begin{definition}[一致收敛]
$\sum u_n(x)$ 在区间 $I$ 上一致收敛于 $S(x)$，若 $\forall \varepsilon>0$，
$\exists N$，当 $n>N$ 时，$\forall x \in I$，$|S_n(x)-S(x)|<\varepsilon$。
\end{definition}

\begin{theorem}[Weierstrass 判别法（M-判别法）]
若 $|u_n(x)| \le M_n$ 对 $\forall x \in I$ 成立且 $\sum M_n$ 收敛，
则 $\sum u_n(x)$ 在 $I$ 上一致收敛。
\end{theorem}

\begin{theorem}[一致收敛的性质]
设 $\sum u_n(x)$ 在 $[a,b]$ 上一致收敛于 $S(x)$，且 $u_n(x)$ 连续：
\begin{itemize}
  \item 和函数 $S(x)$ 在 $[a,b]$ 上连续
  \item 可逐项积分：$\int_a^b S(x)dx = \sum \int_a^b u_n(x)dx$
  \item 若 $\sum u_n'(x)$ 一致收敛，则可逐项求导：
    $S'(x) = \sum u_n'(x)$
\end{itemize}
\end{theorem}

\subsection{幂级数}

\begin{definition}[幂级数]
$\sum_{n=0}^{\infty} a_n (x-x_0)^n$。
\end{definition}

\begin{theorem}[收敛半径]
收敛半径 $R = \frac{1}{\limsup \sqrt[n]{|a_n|}}$（Cauchy-Hadamard 公式）。
若 $\lim |a_{n+1}/a_n|$ 存在，则 $R = \lim |a_n/a_{n+1}|$。
在 $|x-x_0| < R$ 内绝对收敛，在 $|x-x_0| > R$ 外发散，
在 $|x-x_0| = R$ 处需单独判定。
\end{theorem}

\begin{theorem}[幂级数的性质]
在收敛区间内，幂级数的和函数是连续的，且可逐项积分和逐项求导，
所得幂级数有相同的收敛半径。
\end{theorem}

\begin{definition}[Taylor 级数]
若 $f$ 在 $x_0$ 处任意次可导，则其 Taylor 级数为
$\sum_{n=0}^{\infty} \frac{f^{(n)}(x_0)}{n!}(x-x_0)^n$。
当 $x_0=0$ 时称为 Maclaurin 级数。
\end{definition}

\subsection{Fourier 级数}

\begin{definition}[Fourier 级数]
设 $f$ 以 $2\pi$ 为周期且在 $[-\pi,\pi]$ 上可积，其 Fourier 级数为
\[
f(x) \sim \frac{a_0}{2} + \sum_{n=1}^{\infty} (a_n \cos nx + b_n \sin nx),
\]
其中 Fourier 系数：
$a_n = \frac{1}{\pi}\int_{-\pi}^{\pi} f(x) \cos nx dx$，
$b_n = \frac{1}{\pi}\int_{-\pi}^{\pi} f(x) \sin nx dx$。
\end{definition}

\begin{theorem}[Dirichlet 收敛定理]
若 $f$ 以 $2\pi$ 为周期且在 $[-\pi,\pi]$ 上分段光滑，则其 Fourier 级数收敛于
$\frac{f(x^+)+f(x^-)}{2}$。在连续点处收敛于 $f(x)$。
\end{theorem}

\begin{definition}[正弦级数与余弦级数]
奇函数：$a_n=0$，得到正弦级数。偶函数：$b_n=0$，得到余弦级数。
可通过对函数作奇延拓或偶延拓来获得。
\end{definition}

\begin{definition}[Parseval 等式]
$\frac{1}{\pi}\int_{-\pi}^{\pi} f^2(x) dx = \frac{a_0^2}{2} + \sum_{n=1}^{\infty} (a_n^2+b_n^2)$。
\end{definition}

\end{document}
